\sectionkicker{Source-backed technical notes}
\subsection*{Capability \& Evaluation — benchmarkの外で能力をどう確かめるか: Technical Notes}
本欄は記事本文で圧縮した一次資料上の情報を、比較・再検証しやすい形で整理したものである。時系列、確認済みの主張、留意点、一次資料URLを掲載する。
\medskip
\subsection*{Theme at a glance}
\addcontentsline{toc}{subsection}{Theme at a glance}
\begin{center}
\footnotesize
\begin{tabularx}{\linewidth}{@{}p{0.20\linewidth}p{0.10\linewidth}p{0.14\linewidth}X@{}}
\toprule
資料 & 位置づけ & 種別 & 時系列 \\
\midrule
An OpenAI model has disproved a central conjecture in discrete geometry & 主要資料 & 公式情報 & 2026-05-20 (公式公開) \\
A shared playbook for trustworthy third party evaluations & 補足資料 & 評価ガイダンス & 2026-05-29 (公式公開) \\
\bottomrule
\end{tabularx}
\end{center}
\normalsize

\begin{technicalnote}{An OpenAI model has disproved a central conjecture in discrete geometry}{主要資料}
\begingroup
% reader-facing Technical Notes break policy
\widowpenalty=10000
\clubpenalty=10000
\displaywidowpenalty=10000
\begin{tabularx}{\linewidth}{@{}>{\bfseries}p{0.22\linewidth}X@{}}
組織 & OpenAI \\
種別 & 公式情報 \\
時系列 & 2026-05-20 (公式公開) \\
\end{tabularx}
\smallskip
{\bfseries 一次資料から整理したtechnical points}
\begin{itemize}[leftmargin=1.5em,itemsep=0.35em]
\item \textbf{一次情報で確認できる事実}: 保存済み一次資料では「An OpenAI model has disproved a central conjecture in discrete geometry」が2026-05-20に確認できる。
\end{itemize}
\begin{minipage}{\linewidth}
% reader-facing Technical Notes coherent tail group
\begin{itemize}[leftmargin=1.5em,itemsep=0.35em]
\item \textbf{Vendor claim}: OpenAIの公式feedでは、同社modelがdiscrete geometryのunit-distance problemに関する長年のconjectureを反証したと報告している。
\end{itemize}
{\bfseries 読む際の境界}
\begin{itemize}[leftmargin=1.5em,itemsep=0.35em]
\item \textbf{分析上の留意点}: 当該first-party feedは公開時系列とvendorが説明したscopeの確認には使えるが、詳細なcapability・benchmark・safety上の主張は本号で独立再現していない。
\end{itemize}
\begin{samepage}
{\bfseries 一次資料}
\begingroup\sloppy
\begin{itemize}[leftmargin=1.5em,itemsep=0.25em]
\item {\scriptsize\url{https://openai.com/index/model-disproves-discrete-geometry-conjecture}}
\end{itemize}
\endgroup
\end{samepage}
\end{minipage}
\endgroup
\end{technicalnote}
\medskip

\begin{technicalnote}{A shared playbook for trustworthy third party evaluations}{補足資料}
\begingroup
% reader-facing Technical Notes break policy
\widowpenalty=10000
\clubpenalty=10000
\displaywidowpenalty=10000
\begin{tabularx}{\linewidth}{@{}>{\bfseries}p{0.22\linewidth}X@{}}
組織 & OpenAI \\
種別 & 評価ガイダンス \\
時系列 & 2026-05-29 (公式公開) \\
\end{tabularx}
\smallskip
{\bfseries 一次資料から整理したtechnical points}
\begin{itemize}[leftmargin=1.5em,itemsep=0.35em]
\item \textbf{一次情報で確認できる事実}: 保存済み一次資料では「A shared playbook for trustworthy third party evaluations」が2026-05-29に確認できる。
\item \textbf{Vendor claim}: OpenAIの公式feedでは、frontier modelのthird-party evaluationについてcapability、safeguard、evaluation validityを扱うplaybookを公開している。
\end{itemize}
{\bfseries 読む際の境界}
\begin{itemize}[leftmargin=1.5em,itemsep=0.35em]
\item \textbf{分析上の留意点}: 当該first-party feedは公開時系列とvendorが説明したscopeの確認には使えるが、詳細なcapability・benchmark・safety上の主張は本号で独立再現していない。
\end{itemize}
\begin{samepage}
{\bfseries 一次資料}
\begingroup\sloppy
\begin{itemize}[leftmargin=1.5em,itemsep=0.25em]
\item {\scriptsize\url{https://openai.com/index/trustworthy-third-party-evaluations-foundations}}
\end{itemize}
\endgroup
\end{samepage}
\endgroup
\end{technicalnote}
\medskip
